\documentclass[11pt]{article}
\usepackage[margin=1in]{geometry}
\usepackage{amsmath,amssymb,amsfonts,mathtools}
\usepackage{array,booktabs,enumitem,graphicx,xcolor,hyperref}
\usepackage[T1]{fontenc}
\usepackage[utf8]{inputenc}
\title{Self-avoiding effective strings in lattice gauge theories}
\author{Source-backed arXiv sample}
\date{}
\begin{document}
\maketitle
\section{Extracted Lists}

The list environments below are extracted from arXiv LaTeX source and wrapped for pdfLaTeX validation.

\begin{description}
  \item[i)] it is  composed, at large distances, by $D-2$
free bosons, describing the transverse displacements of the string;
  \item[ii)] each bosonic field is compactified on a circle of radius
$R_t$ related to the finite thickness of the flux tube. This allows
fermionizing  the model, which  in turn guarantees that the colour
flux tube does not self-overlap;
  \item[iii)] the boundary conditions on  quark lines are dictated
by the lattice gauge theory (LGT) and fulfil a sort of orthogonality
constraint such that, whenever the quark and antiquark lines coincide,
 the effective string vanishes . This   forces the adimensional radius
$\nu=\sqrt{\pi\sigma}R_t$  to be exactly $\nu=1/4$.
\end{description}
\end{document}
